	
\cnabstract
变电站是电力系统运行维护中的关键节点，其内部存在高电压、强电磁干扰、设备密集、通道狭窄以及动态障碍物等复杂因素。传统人工巡检和辅助作业方式在安全性、实时性和作业效率方面存在一定局限，引入具备自主感知、定位建图和路径规划能力的移动机器人，是提升变电站运维智能化水平的重要途径。围绕变电站运维场景下机器人自主导航的实际需求，本文开展面向变电站运维的搬运机器人导航系统设计与实现研究。

首先，本文分析了变电站运维机器人的功能需求，设计了由传感器模块、中央控制模块、底层运动控制模块和电源模块组成的机器人硬件平台，并结合激光雷达、双目相机等传感器构建环境感知与导航系统框架。在此基础上，介绍了机器人导航所需的坐标系变换、点云聚类等基础理论，为后续环境感知、地图构建和路径规划提供统一的数据表达基础。

其次，针对机器人在复杂动态环境中对障碍物实时感知的需求，本文设计了基于深度图像和点云数据融合的轻量级障碍物检测与跟踪方法。该方法结合U深度图检测器和DBSCAN点云检测器的检测结果，通过包围盒匹配与融合提升障碍物检测的鲁棒性；同时引入基于特征的数据关联、恒加速度运动模型和点云投票策略，对动态障碍物进行状态估计与识别。在地图构建方面，本文采用局部占用栅格更新为规划器提供可更新的环境约束。

最后，面向变电站场景中的复杂障碍分布和实时避障需求，本文设计了全局参考引导与局部滚动优化相结合的分层路径规划方法。全局层采用目标偏向RRT生成二维拓扑可行参考路径，并结合直连检查、线段碰撞检测、超时处理和路径平滑后验校验提高工程可用性；局部层构建学习增强MPC规划器，在承接障碍物边界和动态状态估计结果的基础上，引入点级障碍表示、动态预测点流和学习型距离特征，生成满足运动约束、平滑约束和避障要求的局部轨迹。同时，本文设计MPC直接速度输出接口，将滚动优化得到的首个最优控制量转换为全向底盘可执行的速度指令。

\cnkeyword{变电站运维，移动机器人，自主导航，环境感知，路径规划}
